\documentclass[14pt,a4paper]{ctexart}

\usepackage[top=2.2cm,bottom=2.2cm,left=2.2cm,right=2.2cm,headheight=15pt]{geometry}
\usepackage{amsmath,amssymb}
\usepackage{enumitem}
\usepackage{xcolor}
\usepackage{fancyhdr}

\pagestyle{fancy}
\fancyhf{}
\fancyhead[L]{\textbf{OCR 识别结果汇总}}
\fancyhead[R]{\textbf{ocr-pic-set1}}
\fancyfoot[C]{\thepage}
\renewcommand{\headrulewidth}{0.8pt}

\begin{document}

\begin{center}
{\Huge\bfseries 三张图片 OCR 识别结果}

\bigskip
{\large\color{red} 请逐题对照原图确认，确认后告知哪些题需要解答。}
\end{center}

% ==================== 图片 1 ====================
\section*{图片 1}

\fbox{%
\parbox{\dimexpr\textwidth-2\fboxsep-2\fboxrule\relax}{%
\textbf{题目：}计算
\[
\int_L x\sin y\,dx \;-\; y\sin x\,dy
\]
其中 $L$ 为从点 $A(\pi,0)$ 到点 $B(0,\pi)$ 的直线段。%
}%
}

\bigskip

\color{gray}
OCR 原始输出：\texttt{| zsinyaa 一 ySInZ dy，其中 世 为 从 点 A(ry0) 到 局 刀 (0,r) 的 直线 段 .}
\color{black}

% ==================== 图片 2 ====================
\newpage
\section*{图片 2（教材第十一章习题）}

\fbox{%
\parbox{\dimexpr\textwidth-2\fboxsep-2\fboxrule\relax}{%
\textbf{第3题.} 计算
\[
\oint_L (2xy - x^2)\,dx + (x + y^2)\,dy
\]
$L$ 是由抛物线 $y^2 = x$ 和 $y = x^2$ 所围成的区域的边界曲线，
方向为\textbf{顺时针}方向。

\medskip

\textbf{第4题.} 计算
\[
\int_L (2a - y)\,dx + x\,dy
\]
其中 $L$ 为摆线 $x = a(t - \sin t)$，$y = a(1 - \cos t)$
上由 $t = 0$ 到 $t = 2\pi$ 的一段弧。

\medskip

\textbf{第5题.} 计算
\[
\int_L y\,dx + z\,dy + x\,dz
\]
其中 $L$ 为螺旋线 $x = a\cos t$，$y = a\sin t$，$z = bt$
从 $t = 0$ 到 $t = 2\pi$ 的弧段。

\medskip

\textbf{第6题.} \quad OCR 模糊，仅识别出部分文字：
\small
\[
\int_L (y - z)dx + \cdots\,dz,\qquad
\text{曲线 } \begin{cases} x=? \\ y=? \\ z=? \end{cases}
\text{ 从点 }(0,0,0)\text{ 到点 }(1,1,1)
\]%
}%
}

\bigskip

\color{gray}
OCR 原始输出：\texttt{6. | 一 z)dz 有 de 其 中 卫 为 曲线 | 从 点 (0,0,0) 到 点 (1,1,1) 普 去 灵 的 弧 段}
\color{black}

\begin{center}
\color{red}\large
⚠ 第6题OCR严重模糊，请对照原图补充完整。
\end{center}

% ==================== 图片 3 ====================
\newpage
\section*{图片 3（高等数学练习与提高 三）}

\fbox{%
\parbox{\dimexpr\textwidth-2\fboxsep-2\fboxrule\relax}{%
\textbf{第1题.} 力场由方向为 $y$ 轴的负方向，大小等于作用点的横坐标的平方构成。求当质点沿抛物线 $x = y^2 + 1$ 从点 $(1,0)$ 移动到点 $(2,1)$ 时，力场所做的功。

\medskip

\textbf{第2题.} 将曲线积分 $\displaystyle\int_L e^x(1 - \cos y)dx - e^x(1 - \sin y)dy$ 化为形如
$\displaystyle\int_L P(x,y)dx + Q(x,y)dy$ 的积分，并计算它。其中 $L$ 为上半椭圆
$x^2 + 2y^2 = 2$，从点 $(-1,0)$ 到 $(1,0)$ 的弧段。

\medskip

\textbf{第3题.} 将下列对坐标的曲线积分化为对弧长的曲线积分：

\begin{enumerate}[label=(\arabic*),leftmargin=*]
  \item 计算 $\displaystyle\int_L x^2y\,dx + y\,dy$，$L$ 为曲线 $y = x^2$ 从点 $(0,0)$ 到点 $(1,1)$ 的弧段。
  \item 计算 $\displaystyle\int_L P\,dx + Q\,dy + R\,dz$，$L$ 为曲线 $x = t$，$y = t^2$，$z = t^3$ 上相对应 $t$ 从 $0$ 到 $1$ 的曲线弧段。
\end{enumerate}%
}%
}

\bigskip

\color{gray}
OCR 原始输出（图片3）：\texttt{67 高 等 数学 练习 与 提高 (三 ) ... 有 B 炎 题 ...}
\color{black}

\begin{center}
\color{red}\large
⚠ 图片3 的 OCR 最模糊，第1、2题的被积函数和条件可能不准，\textbf{请务必逐题核对原图}。
\end{center}

% ==================== 汇总 ====================
\newpage
\section*{汇总清单}

\begin{enumerate}[leftmargin=*,itemsep=8pt]
  \item[\textbf{图1-1}] $\displaystyle\int_L x\sin y\,dx - y\sin x\,dy$，
        $L$：从 $A(\pi,0)$ 到 $B(0,\pi)$ 的直线段
  \item[\textbf{图2-3}] $\displaystyle\oint_L (2xy - x^2)dx + (x + y^2)dy$，
        $L$：$y^2=x$ 与 $y=x^2$ 围成区域的边界（顺时针）
  \item[\textbf{图2-4}] $\displaystyle\int_L (2a - y)dx + x\,dy$，
        $L$：摆线 $x=a(t-\sin t), y=a(1-\cos t)$，$t:0\to2\pi$
  \item[\textbf{图2-5}] $\displaystyle\int_L y\,dx + z\,dy + x\,dz$，
        $L$：螺旋线 $x=a\cos t, y=a\sin t, z=bt$，$t:0\to2\pi$
  \item[\textbf{图2-6}] OCR 模糊不清，待确认
  \item[\textbf{图3-1}] 变力沿曲线做功问题，待确认
  \item[\textbf{图3-2}] 曲线积分化简与计算，上半椭圆路径，待确认
  \item[\textbf{图3-3}] 将对坐标的曲线积分化为对弧长的曲线积分（2小题），待确认
\end{enumerate}

\end{document}
